% Fragmento público generado desde propietarios íntegros.
% Residencia: 52.11.2.
% Fuente: ../incoming/radion_monografia_20260827/manuscrito/sections/03_cierre_exacto.tex:286-309
\subsubsection{Structural characterization of \(\pi\) by compatibility of sections}

Reorder \eqref{eq:cierre-unitario} gives
\begin{equation}
\boxed{
2\pih=
\frac{1+\Phi_T-\epsone}
{\frac5{36}+\frac{25}{9}\alphah}.}
\label{eq:pi-radion}
\end{equation}
This equality may be called a new definition of \(\pi\) within the
radion chart if \emph{structural definition} is understood as an
exact characterization by compatibility of sections. It is not a second
independent generator: \(\Delta_4\) and the carry operator already receive the
coinductive section of \(\pi\). The novelty is that the turn is
characterized by the gluing among direct publication, torsion, and memory.

\begin{narrativebox}
The equation \eqref{eq:pi-radion} does not say that we first do not know \(\pi\) and then guess it by closing an angle. It says that the \(\pi\) generated by the nonadic connection is exactly the one that makes two internal readings of the same state compatible. It is an identity of intrinsic recognition, not a metrological calibration.
\end{narrativebox}
